% ---
% titre: Puissances et racines
% theme: NO
% chapitre: Puissances et racines
% sous_chapitre: Calcul réfléchi et mental
% niveau: [9e]
% type: EVAL
% tags: [c-puissance,p-question-TS]
% difficulte: 1
% corrige: true
% ---

\question
Quelle puissance!

\begin{parts}
\vspace{20pt}\part[3] Effectue les calculs suivants.	

\vspace{10pt}
\setlength{\columnsep}{50pt}
\begin{multicols}{3}
\begin{enumerate}[leftmargin=*,align=parleft, label=\alph*), itemsep=15pt]
\item $8^2 =$ \sol{\;\bl}{$64$}
\item $11^0=$ \sol{\;\bl}{$1$}
\item $15^1=$ \sol{\;\bl}{$15$}
\item $2^6=$ \sol{\;\bl}{$64$}
\item $10\,000\,000 = 10^{\sol{\bl[30pt]}{7}}$
\item $\sol{\bl[40pt]}{30}^2= 900$
\end{enumerate}
\end{multicols}
\setlength{\columnsep}{10pt}

\vspace{20pt}\part[2] Complète les phrases suivantes. 

\vspace{10pt}
\begin{enumerate}[leftmargin=*,align=parleft, label=\alph*), itemsep=15pt]
\item Onze puissance deux est égal à \sol{\;\bl}{$121$}.
\item le carré de \sol{\;\bl}{$60$} est égal à trois mille six cent.
\item Cent vingt-cinq est le cube de \sol{\;\bl}{$5$}.
\item Le cube de \sol{\;\bl}{$3$} est égal à vingt-sept.
\end{enumerate}

\begin{solution}[0pt]
\textbf{0.5pt par réponse}
\end{solution}
\end{parts}